\documentclass[11pt]{article}
\usepackage[margin=1in]{geometry}
\usepackage{amsmath, amssymb, amsthm}
\usepackage{mathtools}
\usepackage{hyperref}

\newtheorem{theorem}{Theorem}
\newtheorem{definition}[theorem]{Definition}
\newtheorem{lemma}[theorem]{Lemma}
\newtheorem{remark}[theorem]{Remark}

\title{Random Variables}
\author{math-foundations / node 22}
\date{}

\begin{document}
\maketitle

\subsection*{First, an example}

Before the abstract definition, picture a single fair die roll. Take the
sample space $\Omega = \{1, 2, 3, 4, 5, 6\}$, with $\mathcal{F}$ the power
set and $P$ the uniform measure. Define two functions on $\Omega$:
\[
   X(\omega) = \omega, \qquad
   Y(\omega) = \begin{cases} 1 & \text{if } \omega \text{ is even}, \\ 0 & \text{otherwise}. \end{cases}
\]
\emph{Read aloud:} ``X of omega equals omega; Y of omega equals one if omega is even, zero otherwise.''
Here $X$ takes values in $\{1, \dots, 6\}$ (the number shown), while $Y$
takes values in $\{0, 1\}$ (an \emph{indicator} of the event ``even''). The
event ``$X = 4$'' is the singleton $\{4\} \subset \Omega$, which lies in
$\mathcal{F}$ --- so $X$ qualifies as a random variable. Keep this picture
in mind as we generalize below.

\section{Setting}

Fix a probability space $(\Omega, \mathcal{F}, P)$ and write $\mathcal{B}(\mathbb{R})$ for the Borel
$\sigma$-algebra of $\mathbb{R}$ (the smallest $\sigma$-algebra containing the open intervals).

\begin{definition}[Random variable]
A function $X : \Omega \to \mathbb{R}$ is a \emph{random variable} if
\[
   X^{-1}(B) := \{\omega \in \Omega : X(\omega) \in B\} \;\in\; \mathcal{F}
   \qquad \text{for every } B \in \mathcal{B}(\mathbb{R}).
\]
\emph{Read aloud:} ``the pre-image of B under X --- that is, the set of omega in Omega with X of omega in B --- is an element of script F, for every Borel set B in R.''
\end{definition}

\begin{remark}
Equivalently, $X$ is a random variable iff $\{X \le x\} \in \mathcal{F}$ for every $x \in \mathbb{R}$.
This follows because $\{(-\infty, x] : x \in \mathbb{R}\}$ generates $\mathcal{B}(\mathbb{R})$ and
$X^{-1}$ commutes with countable unions, intersections, and complements.
\end{remark}

\begin{definition}[Distribution function]
The \emph{cumulative distribution function} (CDF) of a random variable $X$ is
\[
   F_X : \mathbb{R} \to [0,1], \qquad F_X(x) := P(X \le x) = P\big(X^{-1}((-\infty, x])\big).
\]
\emph{Read aloud:} ``F sub X maps R to the unit interval; F sub X of x is defined as the probability that X is at most x, equivalently the P-measure of the pre-image of negative-infinity to x under X.''
The function $F_X$ is non-decreasing, right-continuous, with $F_X(-\infty)=0$ and $F_X(+\infty)=1$.
\end{definition}

\section{Algebra of random variables}

We now prove that the sum of two random variables is again a random variable. This is the
\emph{prototype} for showing that products, maxima, minima, and limits of random variables
remain random variables.

\begin{lemma}[Pre-image of an open set]\label{lem:open}
Let $X : \Omega \to \mathbb{R}$ satisfy $\{X \le x\} \in \mathcal{F}$ for all $x \in \mathbb{R}$.
Then for every open set $U \subseteq \mathbb{R}$, $X^{-1}(U) \in \mathcal{F}$.
\end{lemma}

\begin{proof}
Every open subset of $\mathbb{R}$ is a countable union of open intervals. For an interval $(a,b)$,
\[
   X^{-1}((a,b)) = \{X < b\} \cap \{X > a\} = \Big(\bigcup_{n=1}^\infty \{X \le b - 1/n\}\Big) \cap \{X \le a\}^c,
\]
\emph{Read aloud:} ``the pre-image of the open interval (a, b) under X equals the intersection of X-less-than-b and X-greater-than-a, which equals the countable union over n of X-at-most-b-minus-one-over-n, intersected with the complement of X-at-most-a.''
which is a countable Boolean combination of sets in $\mathcal{F}$, hence in $\mathcal{F}$. Taking
countable unions preserves this.
\end{proof}

\begin{theorem}[Sum of random variables is a random variable]
If $X, Y : \Omega \to \mathbb{R}$ are random variables, then $Z := X + Y$ is a random variable.
\end{theorem}

\begin{proof}
By the remark, it suffices to show $\{Z \le c\} \in \mathcal{F}$ for every $c \in \mathbb{R}$.
Equivalently, since $\{Z \le c\} = \{Z < c\}^c \cup \{Z = c\}$ and complements stay in
$\mathcal{F}$, we instead exhibit $\{Z < c\}$ as a measurable set --- this is the cleaner form
and uses Lemma~\ref{lem:open} indirectly.

Observe that $X(\omega) + Y(\omega) < c$ holds iff there exists a rational $q$ with
\[
   X(\omega) < q \quad \text{and} \quad Y(\omega) < c - q.
\]
\emph{Read aloud:} ``X of omega is less than q, and Y of omega is less than c minus q.''
($\Rightarrow$): If $X(\omega) + Y(\omega) < c$, set $\varepsilon = c - X(\omega) - Y(\omega) > 0$;
by density of $\mathbb{Q}$ in $\mathbb{R}$, pick $q \in \mathbb{Q}$ with
$X(\omega) < q < X(\omega) + \varepsilon/2$. Then $q < X(\omega) + \varepsilon/2$ gives
$c - q > c - X(\omega) - \varepsilon/2 = Y(\omega) + \varepsilon/2 > Y(\omega)$.
($\Leftarrow$): If $X(\omega) < q$ and $Y(\omega) < c - q$, add the inequalities.

Hence
\[
   \{Z < c\} \;=\; \bigcup_{q \in \mathbb{Q}} \Big( \{X < q\} \cap \{Y < c - q\} \Big).
\]
\emph{Read aloud:} ``the event Z-less-than-c equals the union over rational q of the intersection of X-less-than-q and Y-less-than-c-minus-q.''
The collection $\mathbb{Q}$ is countable. Each $\{X < q\} = X^{-1}((-\infty, q))$ and each
$\{Y < c-q\} = Y^{-1}((-\infty, c-q))$ is the pre-image of an open set, hence in $\mathcal{F}$
by Lemma~\ref{lem:open}. A countable union of intersections of $\mathcal{F}$-sets is in
$\mathcal{F}$, so $\{Z < c\} \in \mathcal{F}$. Finally
$\{Z \le c\} = \bigcap_{n=1}^\infty \{Z < c + 1/n\} \in \mathcal{F}$, proving the theorem.
\end{proof}

\begin{remark}
The same density-of-rationals trick proves $XY, \max(X,Y), \min(X,Y)$ are random variables.
Iterating gives that every polynomial in finitely many random variables is a random variable,
and combining with Theorem~3 of the README extends the result to limits.
\end{remark}

\end{document}
